\begin{center}
  \large\textbf{ABSTRACT}
\end{center}

\addcontentsline{toc}{chapter}{ABSTRACT}

\begin{center}
  \textbf{\engtatitle{}}
\end{center}

\vspace{1ex}

\begingroup
\setlength{\tabcolsep}{0pt}

\noindent
\begin{tabularx}{\textwidth}{l >{\centering}m{2em} X}
  Student Name / NRP & : & \name{} / \nrp{} \\
  Department         & : & \engdepartment{} \engfacultyshort{} - ITS \\
  Advisor            & : & 1. \advisor{} \\
                     &   & 2. \coadvisor{} \\
\end{tabularx}
\endgroup

\vspace{2ex}

\noindent\textbf{Abstract}

% Ubah paragraf berikut dengan abstrak dari proposal tugas akhir dalam Bahasa Inggris
\emph{\lipsum[1]}

\vspace{2ex}

\noindent\textbf{Keywords:} keyword 1, keyword 2, keyword 3, keyword 4, keyword 5.

